\begin{helper}
Let $n\ge1$ be an integer, let $1<p<\infty$ and let
$A\colon\mathbb{Z}\to\mathbb{C}$.  Write $A_k=A(k)$, let
$\lambda_k=\lambda_k(n,p)$ be as in \noderef{LambdaCoef}, let
$\epsilon_k=\epsilon_k(n,p)$ be as in \noderef{EpsCoef}, and let
$T_k=T_k(n,p,A,\cdot)$ be the $k$-th summand of the series (19) as in
\noderef{Thm2Term}.  Assume the energy hypothesis
\[
  \sum_{m\ge0}(n+m)\,|A_{n+m}|^{2}<\infty .
\]
Then there are real numbers $b$ and $C_A$ such that
\begin{enumerate}
\item[(i)] $0<b$ and $b\le\tfrac12$;
\item[(ii)] $b\,m\le\lambda_{n+m}$ for every integer $m\ge0$;
\item[(iii)] $0\le C_A$ and $|A_{n+m}|\le C_A$ for every integer $m\ge0$;
\item[(iv)] $\rho^{\lambda_{n+m}}\le\bigl(\rho^{b}\bigr)^{m}$ for every
  $\rho$ with $0<\rho<1$ and every $m\ge0$;
\item[(v)] for every integer $M\ge0$ and every $\rho$ with $0<\rho<1$ the
  series
  \[
    \sum_{m\ge0}(n+m)^{M}\,|A_{n+m}|\,\rho^{\lambda_{n+m}}
  \]
  converges;
\item[(vi)] $T_n(\xi)=A_n$ for every $\xi\ne0$;
\item[(vii)] $|T_{n+m}(\xi)|\le2|A_{n+m}|\,|\xi|^{\lambda_{n+m}}$ for every
  integer $m\ge0$ and every $\xi\ne0$;
\item[(viii)] for every $\rho$ with $0<\rho<1$ the series
  $\sum_{m\ge0}2|A_{n+m}|\,\rho^{\lambda_{n+m}}$ converges.
\end{enumerate}

This is the summability engine behind the series (19), and it is the paper's
"we note that $\lim\lambda_k=\infty$.  Therefore, we can perform any finite
number of term-by-term differentiations in formula (19)"
(\texttt{paper.tex:387}) made quantitative.  The point of (ii) is that the rate
is \emph{linear} in $m$ and not merely positive: by (iv) it converts
$\rho^{\lambda_{n+m}}$ into the geometric $(\rho^{b})^{m}$, which is what makes
(v) hold simultaneously for every polynomial weight $M$ and so permits
differentiating the series any finite number of times.  The engine is
\noderef{LambdaLinearLowerBound} alone; the constant lower bound
$\lambda_k\ge(p-1)/(np)$ of (25) contributes nothing here, being constant in
$k$, so it yields neither $\lambda_k\to\infty$ nor any decay.

The numbers $b$ and $C_A$ are produced existentially because no property of
the particular witnesses $\min\bigl(\gamma_n/(n+1),\tfrac12\bigr)$ and
$\bigl(\sum_m(n+m)|A_{n+m}|^{2}\bigr)^{1/2}$ is used downstream --- only the
inequalities listed above.
\end{helper}
\begin{proof}
\emph{Step 0: the four inputs.}  Write $\gamma_n=\gamma_n(p)$ for the exponent
of \noderef{GammaExp}.  We use:
\begin{equation}\tag{1}
  \gamma_n>0 ,
\end{equation}
which is \noderef{GammaRatioMinAtOne} applied at the index $n\ge1$;
\begin{equation}\tag{2}
  \lambda_n=0,\qquad \epsilon_n=0 ,
\end{equation}
which is (23) of \noderef{LambdaEpsEstimates};
\begin{equation}\tag{3}
  |\epsilon_{n+m}|\le1\qquad(m\ge0),
\end{equation}
which for $m=0$ is (2) and for $m\ge1$ is (24) of
\noderef{LambdaEpsEstimates}, since then $n+1\le n+m$; and
\begin{equation}\tag{4}
  \lambda_k\ge\frac{\gamma_n}{n+1}\,k\qquad(k\ge n+1),
\end{equation}
which is \noderef{LambdaLinearLowerBound}.

\emph{Step 1: the rate $b$.}  Put
\[
  b=\min\Bigl(\frac{\gamma_n}{n+1},\ \frac12\Bigr).
\]
Then $b\le\tfrac12$ by definition, and $b>0$ because $\gamma_n/(n+1)>0$ by (1)
and $n+1>0$, and because $\tfrac12>0$; this is (i).

For (ii), fix $m\ge0$.  Since $b\le\gamma_n/(n+1)$ and $m\ge0$,
\[
  b\,m\le\frac{\gamma_n}{n+1}\,m ,
\]
so it suffices to bound the right-hand side by $\lambda_{n+m}$.  If $m=0$ this
right-hand side is $0$ and $\lambda_{n+0}=\lambda_n=0$ by (2), so the two are
equal.  If $m\ge1$ then $n+m\ge n+1$, so (4) applies at $k=n+m$ and gives
$\lambda_{n+m}\ge\frac{\gamma_n}{n+1}(n+m)$; and since
$\frac{\gamma_n}{n+1}\ge0$ and $m\le n+m$ (as $n\ge0$) we have
$\frac{\gamma_n}{n+1}m\le\frac{\gamma_n}{n+1}(n+m)\le\lambda_{n+m}$.  In both
cases $b\,m\le\lambda_{n+m}$.

\emph{Step 2: the coefficient bound $C_A$.}  Let
\[
  S=\sum_{m\ge0}(n+m)|A_{n+m}|^{2},
\]
which is a convergent sum of nonnegative terms by hypothesis, so $S\ge0$; put
$C_A=\sqrt S\ge0$.  Fix $m\ge0$.  Since every term of the sum is nonnegative
and the sum converges, the single term at index $m$ is at most the whole sum:
\[
  (n+m)|A_{n+m}|^{2}\le S .
\]
Also $n+m\ge1$, because $n\ge1$ and $m\ge0$, so
$|A_{n+m}|^{2}\le(n+m)|A_{n+m}|^{2}\le S$.  Taking square roots of the
nonnegative numbers $|A_{n+m}|^{2}\le S$ gives $|A_{n+m}|\le\sqrt S=C_A$.
This is (iii).  (No Cauchy--Schwarz inequality is needed: a single term of a
nonnegative convergent series is bounded by its sum.)

\emph{Step 3: geometric domination, (iv).}  Let $0<\rho<1$ and $m\ge0$.  Since
$\rho>0$, the real-power laws give
$\bigl(\rho^{b}\bigr)^{m}=\rho^{bm}$.  Since $0<\rho\le1$, the function
$t\mapsto\rho^{t}$ is nonincreasing on $\mathbb{R}$, and $bm\le\lambda_{n+m}$
by (ii), so
\[
  \rho^{\lambda_{n+m}}\le\rho^{bm}=\bigl(\rho^{b}\bigr)^{m} .
\]

\emph{Step 4: summability at every polynomial weight, (v).}  Fix an integer
$M\ge0$ and $\rho$ with $0<\rho<1$, and put $q=\rho^{b}$.  Then $q>0$, being a
real power of a positive number, and $q<1$, because $0<\rho<1$ and $b>0$.  By
Step 3,
\begin{equation}\tag{5}
  \rho^{\lambda_{n+m}}\le q^{m}\qquad(m\ge0).
\end{equation}

Next, an elementary inequality: for every real $x\ge0$,
\begin{equation}\tag{6}
  (x+1)^{M}\le2^{M}\bigl(x^{M}+1\bigr).
\end{equation}
Indeed if $x\le1$ then $x+1\le2$, so $(x+1)^{M}\le2^{M}\le2^{M}(x^{M}+1)$
because $x^{M}\ge0$; and if $x>1$ then $x+1\le2x$, so
$(x+1)^{M}\le(2x)^{M}=2^{M}x^{M}\le2^{M}(x^{M}+1)$.  Consequently, for $m\ge0$,
\begin{equation}\tag{7}
  (n+m)^{M}\le\bigl((n+1)(m+1)\bigr)^{M}=(n+1)^{M}(m+1)^{M}
  \le(n+1)^{M}2^{M}\bigl(m^{M}+1\bigr),
\end{equation}
where the first inequality holds because
$(n+1)(m+1)=nm+n+m+1\ge n+m$ and both bases are nonnegative, and the last
step is (6) at $x=m$ multiplied by the nonnegative factor $(n+1)^{M}$.

Combining (7), $|A_{n+m}|\le C_A$ from Step 2, and (5) --- all three factors
being nonnegative --- we get, for every $m\ge0$,
\[
  (n+m)^{M}|A_{n+m}|\,\rho^{\lambda_{n+m}}
  \le(n+1)^{M}2^{M}\bigl(m^{M}+1\bigr)\,C_A\,q^{m}
  =\bigl((n+1)^{M}2^{M}C_A\bigr)\,m^{M}q^{m}
   +\bigl((n+1)^{M}2^{M}C_A\bigr)\,q^{m} .
\]
The majorant on the right is summable in $m$: since $0\le q<1$, the series
$\sum_m m^{M}q^{m}$ converges (the standard polynomial-times-geometric
series) and $\sum_m q^{m}$ converges (geometric series), and a convergent
series stays convergent after multiplication by a constant and addition of
another convergent series.  The terms of the series in (v) are nonnegative,
so by comparison with this majorant that series converges.

\emph{Step 5: the constant summand, (vi).}  Let $\xi\ne0$.  By (2),
$\lambda_n=0$ and $\epsilon_n=0$, so the exponent of the modulus factor is
$\lambda_n+n-n=0$ and the second term of the bracket vanishes.  Hence, by
\noderef{Thm2Term},
\[
  T_n(\xi)=\bigl(A_n\xi^{n}+0\cdot\overline{A_n}\,\bar\xi^{\,n}\bigr)
    \,|\xi|^{0}\,\xi^{-n}
  =A_n\,\xi^{n}\,\xi^{-n}=A_n ,
\]
using $|\xi|^{0}=1$ and $\xi^{n}\xi^{-n}=1$, legitimate because $\xi\ne0$ and
so $\xi^{n}\ne0$.

\emph{Step 6: the pointwise bound, (vii).}  First a bound for a general index.
Let $k\ge0$ be an integer and $\xi\ne0$, and write $r=|\xi|>0$.  Since
$|\overline{A_k}|=|A_k|$, $|\bar\xi|=|\xi|=r$ and $\epsilon_k$ is real,
\begin{equation}\tag{8}
  \bigl|A_k\xi^{k}+\epsilon_k\overline{A_k}\,\bar\xi^{\,k}\bigr|
  \le|A_k|\,r^{k}+|\epsilon_k|\,|A_k|\,r^{k}
  =\bigl(1+|\epsilon_k|\bigr)|A_k|\,r^{k}
\end{equation}
by the triangle inequality and multiplicativity of the modulus.  Taking moduli
in the definition of $T_k$ and using $\bigl||\xi|^{a}\bigr|=r^{a}$ for real $a$
(the modulus factor is a positive real) and $|\xi^{-n}|=r^{-n}$,
\[
  |T_k(\xi)|
  =\bigl|A_k\xi^{k}+\epsilon_k\overline{A_k}\bar\xi^{\,k}\bigr|
     \cdot r^{\lambda_k+n-k}\cdot r^{-n}
  \le\bigl(1+|\epsilon_k|\bigr)|A_k|\,
     r^{k}\,r^{\lambda_k+n-k}\,r^{-n}
  =\bigl(1+|\epsilon_k|\bigr)|A_k|\,r^{\lambda_k},
\]
where the inequality is (8) multiplied by the nonnegative factors
$r^{\lambda_k+n-k}$ and $r^{-n}$, and the last step adds the exponents of the
positive base $r$:
$k+(\lambda_k+n-k)+(-n)=\lambda_k$.  Now take $k=n+m$.  By (3) we have
$1+|\epsilon_{n+m}|\le2$, and $|A_{n+m}|\ge0$ and $r^{\lambda_{n+m}}\ge0$, so
\[
  |T_{n+m}(\xi)|\le\bigl(1+|\epsilon_{n+m}|\bigr)|A_{n+m}|\,r^{\lambda_{n+m}}
  \le2|A_{n+m}|\,|\xi|^{\lambda_{n+m}} ,
\]
which is (vii).

\emph{Step 7: (viii).}  Apply (v) with $M=0$: since $(n+m)^{0}=1$, the series
$\sum_m|A_{n+m}|\rho^{\lambda_{n+m}}$ converges for every $\rho\in(0,1)$.
Multiplying a convergent series of reals by the constant $2$ leaves it
convergent, and $2|A_{n+m}|\rho^{\lambda_{n+m}}
=2\bigl(|A_{n+m}|\rho^{\lambda_{n+m}}\bigr)$ termwise, which gives (viii).

With $b$ from Step 1 and $C_A$ from Step 2, all eight assertions have been
established, so the existential statement holds.
\end{proof}
